% ICLR 2027 — closed-loop quality-triage system (anonymized).  Skeleton; section order LOCKED to
% STORY_FRAMEWORK.md §故事弧. Do NOT reorder §3-§7 (anti-drift rule 6).
% Replace \documentclass with the official iclr2027 style when available.
\documentclass{article}
\usepackage[utf8]{inputenc}
% ICLR 2027 preamble — packages, theorem environments, anonymized macros.
% Model-name macros are defined to ANONYMIZED strings so the literal internal
% names never appear in the .tex (passes the redline hook + one-line de-anon toggle).

\usepackage{amsmath,amssymb,amsthm,mathtools}
\usepackage{bm}
\usepackage{graphicx}
\usepackage{booktabs}
\usepackage{multirow}
\usepackage{xcolor}
\usepackage{hyperref}
\usepackage{cleveref}

% ── Theorem environments ───────────────────────────────────────────────────────
\theoremstyle{plain}
\newtheorem{theorem}{Theorem}
\newtheorem{proposition}{Proposition}
\newtheorem{lemma}{Lemma}
\newtheorem{corollary}{Corollary}
\theoremstyle{definition}
\newtheorem{assumption}{Assumption}
\newtheorem{remark}{Remark}

% ── Anonymized model-name macros (de-anon = edit RHS only) ─────────────────────
\newcommand{\qvib}{QC-VIB}            % quality-conditional VIB
\newcommand{\visiscore}{5-head IQA}  % 5-dim quality scorer
\newcommand{\visienhance}{QP-Enhance}% quality-preserving enhancement
\newcommand{\qcts}{QC-TS}            % quality-conditional temperature scaling (prior submission)
\newcommand{\itb}{QSB}              % quality-stratified benchmark

% ── Math operators ─────────────────────────────────────────────────────────────
\DeclareMathOperator{\KL}{KL}
\DeclareMathOperator*{\E}{\mathbb{E}}
\DeclareMathOperator{\Var}{Var}
\DeclareMathOperator{\softmax}{softmax}
\DeclareMathOperator{\sigmoid}{sigmoid}
\DeclareMathOperator{\tr}{tr}
\DeclareMathOperator*{\argmax}{arg\,max}
\newcommand{\TV}{\mathrm{TV}}
\newcommand{\Hent}{H}                 % Shannon entropy (nats)
\newcommand{\qbar}{\bar{q}}
\newcommand{\Tw}{T_\omega}            % enhancement map
\newcommand{\Zenh}{Z_{\mathrm{enh}}}
\newcommand{\Zref}{Z_{\mathrm{ref}}}
\newcommand{\Lq}{L_{q_\theta}}        % classifier Lipschitz constant
\newcommand{\Ldp}{\mathcal{L}_{\mathrm{DP}}}
\newcommand{\Ragent}{\mathcal{R}_{\mathrm{agent}}}
\newcommand{\Rdirect}{\mathcal{R}_{\mathrm{direct}}}
\newcommand{\Renh}{\mathcal{R}_{\mathrm{enh}}}
\newcommand{\tauenh}{\tau_{\mathrm{enh}}}
\newcommand{\tauhigh}{\tau_{\mathrm{high}}}
\newcommand{\taulow}{\tau_{\mathrm{low}}}


\title{Closed-Loop Quality-Triage for Reliable Skin-Lesion Diagnosis\\
on Quality-Varying Consumer Imagery}
\author{Anonymous Authors}
\date{}

\begin{document}
\maketitle

\begin{abstract}
% LOCKED hook (anti-drift rule 1): first sentence MUST stay on the closed-loop
% quality-triage agent. Do not reduce to a calibration-only story (rule 7).
% TODO(M3): fill — detect (\visiscore) -> enhance (\visienhance) OR query -> diagnose (\qvib),
% with a 5-theorem end-to-end closure and an expected-risk bound (Theorem~\ref{thm:agent}).
\end{abstract}

% ═══════════════════════════════════════════════════════════════════════════════
\section{Introduction}\label{sec:intro}
% §1.1 problem: consumer-device medical AI, quality varies, passive calibration insufficient
% §1.2 gaps: post-hoc calibration (quality-oblivious) / ensemble & dropout UQ (quality-oblivious) /
%            generic enhancement (no diagnostic-preservation guarantee)
% §1.3 ★ closed-loop quality-triage hook ★
% §1.4 Contributions C1-C4
\paragraph{Contributions.}
% C1: \qvib — first quality-conditional information bottleneck (Prop 1, Lemma 1, Thm 1, Prop 2).
% C2: \visienhance + DP-Loss — diagnosis-preserving enhancement with a mutual-information guarantee
%      (Prop 3, Lemma 3).
% C3: closed-loop agent with an expected-risk bound (Thm 2) + composite-calibration bound (Cor 1).
% C4: \itb benchmark + 25-lever evaluation.

% ═══════════════════════════════════════════════════════════════════════════════
\section{Related Work}\label{sec:related}
% §2.1 Calibration; §2.2 Information Bottleneck; §2.3 Medical Image Enhancement;
% §2.4 IQA + Medical AI Benchmarks.

% ═══════════════════════════════════════════════════════════════════════════════
\section{Quality Assessment and the Quality-Conditional Bottleneck}\label{sec:qvib}
% §3.1 \visiscore (5-dim quality vector, PLCC 0.924 / SRCC 0.895)
% §3.2 \qvib ELBO (Proposition 1)
% §3.3 Quality-Adaptive Prior (Lemma 1, sigmoid parameterisation)
% §3.4 Quality Tokenizer + Attention Drift Bound (Theorem 1)
% §3.5 Quality-Calibrated Entropy Guarantee (Proposition 2)
% Full proofs: Appendix A1.

% ═══════════════════════════════════════════════════════════════════════════════
\section{Diagnosis-Preserving Enhancement}\label{sec:enhance}
% §4.1 Degradation taxonomy (q_3 completeness irreversible)
% §4.2 NAFNet + FiLM architecture
% §4.3 DP-Loss definition
\subsection{Entropy Reduction and Mutual-Information Preservation}\label{sec:prop3lemma3}
% Compact statements; full proofs in Appendix A2.1 / A2.2.
% Compact statements of Prop 3 + Lemma 3 for the main body (§4.4).
% Full proofs: Appendix A2.1 / A2.2 (appendix/A2_prop3_lemma3.tex).

\begin{proposition}[Enhancement Entropy Reduction]\label{prop:entropy}
Under Assumptions A1--A4 (quality improvement, calibrated scoring, monotone prior,
diagnostic preservation), for $x$ in the salvage band and $q'=f_\eta(\Tw(x,q))$,
\begin{equation}
\E_{x}\!\big[\Hent(\hat{p}_{\Tw(x,q)})\big] \;\leq\;
\E_{x}\!\big[\Hent(\hat{p}_x)\big]
- \tfrac{d}{2}\log\tfrac{\sigma^2(\qbar)}{\sigma^2(\qbar')} + 2\beta\sqrt{\epsilon},
\label{eq:prop3}
\end{equation}
where the first negative term is the \qvib{} entropy gap (Lemma~1, strictly positive
since $\qbar'>\qbar \Rightarrow \sigma^2$ shrinks) and the residual $2\beta\sqrt{\epsilon}$
is the diagnostic-preservation slack (Lemma~\ref{lem:dp}).
\end{proposition}

\begin{lemma}[Diagnosis-Preserving Mutual-Information Bound]\label{lem:dp}
Let $q_\theta$ be $\Lq$-Lipschitz in total variation and let the label entropy be
bounded by $M=\log K$. If $\Ldp := \KL(p_\phi^{\mathrm{enh}} \,\|\, p_\phi^{\mathrm{ref}}) \leq \epsilon$, then
\begin{equation}
I(\Zenh; Y) \;\geq\; I(\Zref; Y) - \beta\sqrt{\epsilon},
\qquad \beta = \frac{M\,\Lq}{\sqrt{2}}.
\label{eq:lemma3}
\end{equation}
The $\sqrt{\epsilon}$ rate is Pinsker-optimal; a linear-$\epsilon$ bound would require a
stronger $\chi^2$-divergence assumption (Appendix~A2.2.5). The proof combines Pinsker's
inequality with a Lipschitz coupling argument (Appendix~A2.2).
\end{lemma}

% §4.5 Three-stage training protocol (Stage 1 pixel / Stage 2 DP-Loss / Stage 3 quality hinge)

% ═══════════════════════════════════════════════════════════════════════════════
\section{Closed-Loop Agent}\label{sec:agent}
% §5.1 Dual-channel decision logic (4 channels: direct / cautioned / enhance / query-refuse)
% §5.2 Theorem 2 (expected risk bound) — compact; full proof Appendix A2.3
% §5.3 Corollary 1 (\qvib + \qcts composite ECE bound) — link to prior submission
% §5.4 Agent implementation (ReAct + small LLM + rule fallback)

% ═══════════════════════════════════════════════════════════════════════════════
\section{The \itb{} Benchmark}\label{sec:benchmark}
% §6.1 4 subsets (LQ / HQ / Edge / Diverse); §6.2 quality-stratified protocol; §6.3 9 baselines

% ═══════════════════════════════════════════════════════════════════════════════
\section{Experiments}\label{sec:exp}
% §7.1 Setup; §7.2 Main results (Table 1); §7.3 ★ diagnostic preservation ★;
% §7.4 ★ salvage rate ★; §7.5 universality; §7.6 cross-domain; §7.7 cost-benefit + DCA; §7.8 fairness

% ═══════════════════════════════════════════════════════════════════════════════
\section{Discussion and Limitations}\label{sec:discussion}
% §8.1 connection to prior post-hoc calibration; §8.2 why generative enhancement fails here;
% §8.3 failure-mode taxonomy (3-mode); §8.4 limitations (theory is sketch, synthetic benchmark, etc.)

% ═══════════════════════════════════════════════════════════════════════════════
\section{Conclusion}\label{sec:conclusion}

\appendix
% Appendix A2.1 / A2.2 — full proofs of Proposition 3 and Lemma 3.
% Source-of-truth: project/plans/Prop3_Lemma3_visienhance_theory.md (audited 2026-05-27,
% canonical form: √ε scaling, β = M L_q / √2). Numerical toy: tests/test_theorems_numerical.py.

\section{Proofs for Diagnosis-Preserving Enhancement}\label{app:a2}

\subsection{Assumptions}\label{app:a2-assump}
\begin{assumption}[Quality improvement]\label{as:a1}
$\Tw$ satisfies $\qbar(\Tw(x,q)) > \qbar(x)$ almost surely on the salvage band
$\qbar\in[\tauenh,\tauhigh]$ (outside this band no improvement is guaranteed).
\end{assumption}
\begin{assumption}[Calibrated quality scoring]\label{as:a2}
The frozen scorer $f_\eta:\mathcal{X}\to[0,1]^5$ yields $\qbar$ monotone in ground-truth
degradation severity on paired data (PLCC $0.924$, SRCC $0.895$).
\end{assumption}
\begin{assumption}[Monotone prior]\label{as:a3}
The \qvib{} prior variance $\sigma^2(\qbar)$ is strictly decreasing in $\qbar$ on $(0,1)$
(Lemma~1), i.e.\ $\tfrac{d\sigma^2}{d\qbar} < 0$.
\end{assumption}
\begin{assumption}[Diagnostic preservation]\label{as:a4}
$\Tw$ is trained with a DP loss $\Ldp\leq\epsilon$ (Lemma~\ref{lem:dp-full}), so
$p_\phi(z\mid x_{\mathrm{enh}},q')$ and $p_\phi(z\mid x_{\mathrm{ref}},q_{\mathrm{ref}})$
are $\epsilon$-close in KL.
\end{assumption}

% ═══════════════════════════════════════════════════════════════════════════════
\subsection{Proposition 3 (Variance-Driven Entropy Term, Sketch)}\label{app:a2-prop3}

\paragraph{Scope and empirical status.}
We sketch the directional effect of the quality-adaptive prior variance on predictive
entropy under enhancement. We emphasise at the outset that the \emph{net} entropy of the
predictive distribution is \textbf{not} empirically supported to decrease under
enhancement: a quality-conditioned predictive-entropy probe finds that the mean entropy
\emph{increases} after enhancement on the mixed-severity test split
($\bar{\Hent}_{\mathrm{deg}}=0.1487\!\to\!\bar{\Hent}_{\mathrm{enh}}=0.1572$; two
independent probes agree). The derivation below therefore isolates only the
$\sigma^2$-driven term $-\tfrac{d}{2}\log\tfrac{\sigma^2(\qbar)}{\sigma^2(\qbar')}$ and
its direction; it does \emph{not} establish a net entropy reduction, because the
residual and drift terms are not controlled tightly enough to dominate empirically.
Consequently, the role that Proposition~3 plays in the paper --- linking enhancement to
preserved discriminative information --- is carried not by a net-entropy claim but by
(i)~Lemma~\ref{lem:dp-full}, whose mutual-information lower bound is empirically validated
(experiment E7: $\Delta\mathrm{AUC}_{\mathrm{enh}}=+0.0205$, $95\%$ CI $[+0.005,+0.035]$,
McNemar $p=2.3\times10^{-45}$), and (ii)~the Stage-1 PSNR$\,\geq 30$\,dB non-emptiness
condition on the salvage band (Assumption~\ref{as:a1}). See the Limitations discussion in
the main text, which states that these arguments are sketches rather than fully watertight
proofs.

\begin{proposition}[$\sigma^2$-driven entropy term, directional]\label{prop:entropy-full}\label{prop:entropy}
Under Assumptions~\ref{as:a1}--\ref{as:a4}, for $x$ drawn from the salvage band and
$q'=f_\eta(\Tw(x,q))$, the variance-driven component of the expected predictive entropy
satisfies the directional bound
\begin{equation}
\E_{x\sim p(x\mid\qbar)}\!\big[\Hent(\hat{p}_{\Tw(x,q)})\big] \;\leq\;
\E_{x\sim p(x\mid\qbar)}\!\big[\Hent(\hat{p}_x)\big]
- \tfrac{d}{2}\log\tfrac{\sigma^2(\qbar)}{\sigma^2(\qbar')} + 2\beta\sqrt{\epsilon},
\label{eq:p3-full}
\end{equation}
\emph{provided} the unmodelled mean-drift contribution $\Delta\Hent_\mu$ (Step~4) is
non-positive in expectation. We stress that this proviso is \textbf{not} met empirically
(see the scope note above), so \eqref{eq:p3-full} should be read as characterising the
\emph{sign} of the $\sigma^2$-driven term, not as a certificate of net entropy reduction.
\end{proposition}

\begin{proof}[Derivation (sketch, under Assumptions~\ref{as:a1}--\ref{as:a4}; see the
main-text Limitations for non-watertight regimes)]
We sketch the argument in five steps; it isolates the direction of the
$\sigma^2$-driven term rather than proving a net reduction.

\paragraph{Step 1 (ELBO decomposition of predictive entropy).}
By Jensen's inequality and the \qvib{} ELBO (Proposition~1),
\begin{equation}
\Hent(\hat{p}_x) = \Hent\!\big(\E_{z\sim p_\phi(z\mid x,q)}[q_\theta(y\mid z)]\big)
\;\geq\; \E_{z\sim p_\phi}\!\big[\Hent(q_\theta(y\mid z))\big] - I(Z;Y\mid X=x).
\label{eq:p3-1}
\end{equation}
The Markov chain $X\to Z\to Y$ keeps $I(Z;Y\mid X=x)$ small, so the dominant term is
$\E_z[\Hent(q_\theta(y\mid z))]$.

\paragraph{Step 2 (Bound conditional entropy by encoder variance).}
A second-order Taylor expansion of $\Hent(q_\theta(y\mid z))$ around the posterior mean
$\mu_\phi(x,q)$ gives, with $\Sigma_\phi=\mathrm{diag}(\sigma_\phi^2(x,q))$,
\begin{equation}
\E_z[\Hent(q_\theta(y\mid z))] = \Hent(q_\theta(y\mid\mu_\phi))
+ \tfrac{1}{2}\sum_{j=1}^{d}\sigma_{\phi,j}^2\,\partial^2_{z_j}\Hent + O(\|\sigma_\phi\|^4).
\label{eq:p3-2}
\end{equation}

\paragraph{Step 3 (KL forces $\sigma_\phi^2$ to track $\sigma^2(\qbar)$).}
The \qvib{} KL regulariser is
$\KL(p_\phi(z\mid x,q)\,\|\,r_\psi(z\mid q))
= \tfrac{1}{2}\sum_j[\tfrac{\mu_{\phi,j}^2+\sigma_{\phi,j}^2}{\sigma^2(\qbar)}
- 1 - \log\tfrac{\sigma_{\phi,j}^2}{\sigma^2(\qbar)}]$.
At training equilibrium the first-order condition in $\sigma_{\phi,j}^2$ yields
$\sigma_{\phi,j}^2 = \sigma^2(\qbar)\,(1+\mathcal{O}(\beta^{-1}))$. Substituting into
\eqref{eq:p3-2},
\begin{equation}
\E_z[\Hent(q_\theta(y\mid z))] \approx \Hent(q_\theta(y\mid\mu_\phi))
+ \tfrac{d}{2}\sigma^2(\qbar)\,c(x,q) + O(\sigma^4),
\label{eq:p3-3}
\end{equation}
with $c(x,q)=\tfrac{1}{d}\sum_j\partial^2_{z_j}\Hent|_{z=\mu_\phi}\geq 0$ on average.

\paragraph{Step 4 (Compare $\qbar$ versus $\qbar'$).}
For the enhanced input $x_{\mathrm{enh}}=\Tw(x,q)$ with $\qbar'>\qbar$ (Assumptions~\ref{as:a1},
\ref{as:a2}),
\begin{equation}
\E_z[\Hent(q_\theta(y\mid z_{\mathrm{enh}}))] - \E_z[\Hent(q_\theta(y\mid z_x))]
\approx \tfrac{d}{2}\big(\sigma^2(\qbar') - \sigma^2(\qbar)\big)\,\bar{c} + \Delta\Hent_\mu,
\label{eq:p3-4}
\end{equation}
where $\bar{c}\geq 0$ and $\sigma^2(\qbar')<\sigma^2(\qbar)$ by Assumption~\ref{as:a3} makes the
first term strictly negative. The mean drift $\Delta\Hent_\mu$ is \emph{bounded in magnitude}
by Assumption~\ref{as:a4}: Lemma~\ref{lem:dp-full} gives
$\|p_\phi(z\mid x_{\mathrm{enh}},q') - p_\phi(z\mid x_{\mathrm{ref}},q_{\mathrm{ref}})\|_{\TV}
\leq \sqrt{\epsilon/2}$, and Fannes-style continuity with Lipschitz $q_\theta$ yields
$\|\Delta\Hent_\mu\| \leq 2\beta\sqrt{\epsilon}$ (same Pinsker constant as Lemma~\ref{lem:dp-full}).
We caution that this controls only the \emph{magnitude} of $\Delta\Hent_\mu$, not its sign:
the bound permits $\Delta\Hent_\mu$ to be positive, in which case the mean-drift contribution
can offset or exceed the negative variance-driven term. This is precisely the regime our
measurements indicate (net entropy increases under enhancement), so we do not claim that the
right-hand side of \eqref{eq:p3-4} is negative in practice.

\paragraph{Step 5 (Refine via Lemma~1 monotonicity).}
With $\sigma^2(\qbar)=\sigma_{\min}^2+(\sigma_{\max}^2-\sigma_{\min}^2)(1-\sigmoid(s(\qbar-\qbar_0)))$,
for $\qbar'=\qbar+\Delta\qbar$ with $\Delta\qbar\in[0.1,0.25]$,
\begin{equation}
\log\tfrac{\sigma^2(\qbar)}{\sigma^2(\qbar')}
\geq s\,\Delta\qbar\,(1-\sigmoid(\cdot))\,\tfrac{\sigma_{\max}^2-\sigma_{\min}^2}{\sigma_{\max}^2}
= \Theta(s\,\Delta\qbar) > 0.
\label{eq:p3-5}
\end{equation}
Combining \eqref{eq:p3-4}--\eqref{eq:p3-5} and taking the expectation over $x$ in the
salvage band gives the directional statement \eqref{eq:p3-full}, under the proviso (stated in
the proposition) that $\Delta\Hent_\mu$ is non-positive in expectation. Since that proviso is
not met empirically, we treat \eqref{eq:p3-full} as a characterisation of the sign of the
$\sigma^2$-driven term and defer the operative information-preservation guarantee to
Lemma~\ref{lem:dp-full} (validated by experiment E7) and the PSNR$\,\geq 30$\,dB
non-emptiness condition.
\end{proof}

\paragraph{Magnitude of the $\sigma^2$-driven term.}
For $d=128$, $\sigma_{\max}^2/\sigma_{\min}^2\approx 4$, $s\approx 6$, at $\qbar=0.45\to 0.60$,
the variance-driven term has magnitude
$\tfrac{d}{2}\log\tfrac{\sigma^2(0.45)}{\sigma^2(0.60)} \approx 64\log(1.5) \approx 26$ nats,
whereas the magnitude bound on the controlled residual is
$2\beta\sqrt{\epsilon}\approx 2\cdot0.735\cdot\sqrt{0.05}\approx 0.33$ nats
($\beta=M\Lq/\sqrt{2}\approx 0.735$ for $K=2$, $\epsilon=0.05$).
% VERIFIED (会话35 theory reviewer): beta=M*Lq/sqrt(2)=0.6931*1.5/sqrt(2)=0.735 with M=log2 exact; Lq=1.5 is the spectral-clip design bound. Constants confirmed self-consistent; no numerical change.
We caution against reading this comparison as a non-vacuity certificate: it bounds only the
\emph{modelled} variance-driven and Pinsker-residual terms, and it does \emph{not} account
for the unmodelled mean-drift contribution $\Delta\Hent_\mu$ of Step~4, whose sign is
unconstrained. Empirically the net entropy increases under enhancement, indicating that
$\Delta\Hent_\mu$ dominates in our regime; we therefore do not assert that the net bound
\eqref{eq:p3-full} is realised.

\paragraph{Failure mode.}
If Assumption~\ref{as:a1} fails ($\Tw$ collapses to identity, $\qbar'\approx\qbar$), the
RHS of \eqref{eq:p3-5} vanishes and \eqref{eq:p3-full} reduces to the trivial bound. This is
detected empirically by $\qbar_{\mathrm{enh}}-\qbar_{\mathrm{low}}<0.05$ on the validation set
(the Stage-1 PSNR $\geq 30$\,dB gate).

% ═══════════════════════════════════════════════════════════════════════════════
\subsection{Lemma 3 (DP-Loss $\Rightarrow$ Mutual-Information Lower Bound)}\label{app:a2-lemma3}

\begin{lemma}[restated]\label{lem:dp-full}
Assume (L1) $\Hent(Y)\leq M=\log K$ and (L2) $q_\theta(y\mid z)$ is $\Lq$-Lipschitz in $z$
under total variation. If $\Ldp := \KL(p_\phi^{\mathrm{enh}}\,\|\,p_\phi^{\mathrm{ref}})\leq\epsilon$, then
\begin{equation}
I(\Zenh;Y) \;\geq\; I(\Zref;Y) - \beta\sqrt{\epsilon},
\qquad \beta := \frac{M\,\Lq}{\sqrt{2}}.
\label{eq:l3-full}
\end{equation}
\end{lemma}

\begin{proof}[Derivation (sketch, under the stated assumptions (L1)--(L2); see the main-text
Limitations for non-watertight regimes)]
The argument proceeds in four steps. This is a derivation under the stated assumptions rather
than a fully watertight proof; its empirical instantiation is reported as experiment E7.

\paragraph{Step 1 (Pinsker: KL $\to$ TV).}
By Pinsker's inequality,
\begin{equation}
\delta_{\TV}(p_\phi^{\mathrm{enh}},p_\phi^{\mathrm{ref}})
= \tfrac{1}{2}\|p_\phi^{\mathrm{enh}}-p_\phi^{\mathrm{ref}}\|_1
\leq \sqrt{\tfrac{1}{2}\KL(p_\phi^{\mathrm{enh}}\,\|\,p_\phi^{\mathrm{ref}})} \leq \sqrt{\epsilon/2}.
\label{eq:l3-1}
\end{equation}
(KL is in nats throughout.)

\paragraph{Step 2 (Push TV through the classifier by coupling).}
For the optimal coupling $(\Zenh,\Zref)$, $\mathbb{P}[\Zenh\neq\Zref]=\delta_{\TV}$. With the
$\Lq$-Lipschitz classifier and writing $\bar{\sigma}:=\E\|\Zenh-\Zref\|_2/\delta_{\TV}$ for the
effective latent radius (of order $\sigma(\qbar)$, hence $\qbar$-dependent),
\begin{equation}
\|\hat{p}_{\mathrm{enh}}-\hat{p}_{\mathrm{ref}}\|_{\TV}
= \big\|\E[q_\theta(\cdot\mid\Zenh)-q_\theta(\cdot\mid\Zref)]\big\|_{\TV}
\leq \Lq\,\E\|\Zenh-\Zref\|_2 \leq \Lq\,\bar{\sigma}\,\sqrt{\epsilon/2}.
\label{eq:l3-2}
\end{equation}
We keep the $\bar{\sigma}$ factor explicit rather than folding it into $\Lq$, because
$\bar{\sigma}$ depends on $\qbar$ and absorbing it would silently entangle the reported
constant $\Lq\approx 1.5$ (and hence $\beta=M\Lq/\sqrt{2}$) with the quality level.
% VERIFIED (会话35 theory reviewer): under bar-sigma=Theta(1) the downstream constants are unchanged; beta=M*Lq/sqrt(2)=0.6931*1.5/sqrt(2)=0.735 (M=log2 exact), Lq=1.5 is the spectral-clip design bound. sigma-bar now carried through eq l3-3 and absorbed explicitly. No numerical change.
In the normalised-head regime of our setup the effective radius is itself $O(1)$ after the
spectral-norm clipping that defines $\Lq$, so the subsequent constants are stated for
$\bar{\sigma}=\Theta(1)$; the $\qbar$-dependence is what we make explicit here.

\paragraph{Step 3 (MI gap via continuity of entropy).}
Writing $I(Z;Y)=\Hent(Y)-\Hent(Y\mid Z)$, the drop equals
$\Hent(Y\mid\Zenh)-\Hent(Y\mid\Zref)$. By the Fannes--Audenaert inequality and
\eqref{eq:l3-2},
\begin{equation}
|\Hent(Y\mid\Zenh)-\Hent(Y\mid\Zref)|
\leq M\,\delta_{\TV} + h(\delta_{\TV})
\leq M\,\Lq\,\bar{\sigma}\sqrt{\epsilon/2} + O\!\big(\sqrt{\epsilon}\log(1/\sqrt{\epsilon})\big)
\;\stackrel{\bar{\sigma}=\Theta(1)}{=}\; M\,\Lq\sqrt{\epsilon/2} + O\!\big(\sqrt{\epsilon}\log(1/\sqrt{\epsilon})\big),
\label{eq:l3-3}
\end{equation}
where $h$ is the binary entropy, dominated by its linear term for small $\delta_{\TV}$.

\paragraph{Step 4 (Final bound).}
Taking the worst-case direction,
\begin{equation}
I(\Zref;Y)-I(\Zenh;Y)
\leq M\,\Lq\sqrt{\epsilon/2}
= \frac{M\,\Lq}{\sqrt{2}}\,\sqrt{\epsilon} = \beta\sqrt{\epsilon},
\label{eq:l3-4}
\end{equation}
which is \eqref{eq:l3-full}. The $\sqrt{\epsilon}$ rate is the standard Pinsker scaling for a
KL budget; a linear-$\epsilon$ rate would require a stronger $\chi^2$-divergence budget
($\chi^2\leq\epsilon$ in place of KL; Appendix~A2.2.5). We do not claim this rate is optimal,
only that it is the rate Pinsker's inequality delivers under the KL assumption (L2).
\end{proof}

\paragraph{Tight constants.}
$M=\log 2\approx 0.693$ ($K=2$); the normalised \qvib{} head has $\Lq\approx 1.5$
(spectral-norm clipped); $\epsilon=0.05$ (Stage-2 target). Then
$\beta\sqrt{\epsilon}\approx \tfrac{0.693\cdot 1.5}{\sqrt{2}}\cdot\sqrt{0.05}\approx 0.165$ nats
($K=2$), or $\approx 0.46$ nats for $K=7$ --- well within the inductive gap from the \qvib{}
ELBO ($I(Z;Y)\approx 0.5$--$0.6$ nats on the benchmark). A numerical toy
(paired Gaussian latents + Lipschitz linear classifier) confirms the regression slope of
the MI drop against $\sqrt{\epsilon}$ is finite and below $\beta_{\mathrm{theory}}\approx 0.735$.

% % Appendix A1 — full proofs for the Quality-Conditional Bottleneck.
% Source-of-truth: project/plans/V-QIB数学推导.md (Claude 4.7 Opus, 2026-05-06).
% Results: Proposition 1 (ELBO), Lemma 1 (prior monotonicity),
% Lemma 2 (softmax perturbation), Theorem 1 (attention drift), Proposition 2 (entropy monotonicity).
% Anonymized: model names via macros (\qvib, \visiscore). Never write the literal internal names.

\section{Proofs for the Quality-Conditional Bottleneck}\label{app:a1}

\subsection{Setup and Notation}\label{app:a1-setup}
Let $\mathcal{X}\subset\mathbb{R}^{H\times W\times 3}$ be the input space, $\mathcal{Y}=\{1,\dots,K\}$ the
diagnostic labels, and $\mathcal{Z}=\mathbb{R}^d$ the latent space. The frozen scorer
$\visiscore$ maps an image to a quality vector $q\in[0,1]^5$ with mean
$\qbar=\tfrac{1}{5}\sum_{d=1}^5 q_d$. We write $p_\phi(z\mid x,q)$ for the stochastic encoder,
$q_\theta(y\mid z)$ for the classifier, and $r_\psi(z\mid q)$ for the quality-adaptive marginal
prior. All KL divergences are in nats; densities are with respect to Lebesgue measure.

% ═══════════════════════════════════════════════════════════════════════════════
\subsection{Proposition 1 (\qvib{} Evidence Lower Bound)}\label{app:a1-prop1}

\begin{proposition}[restated]\label{prop:elbo-full}
For any encoder $p_\phi(z\mid x,q)$, classifier $q_\theta(y\mid z)$, quality-adaptive prior
$r_\psi(z\mid q)$, and $\beta>0$, the conditional log-likelihood obeys
\begin{equation}
\log p(y\mid x,q) \;\geq\;
\E_{z\sim p_\phi(z\mid x,q)}\!\big[\log q_\theta(y\mid z)\big]
- \beta\,\KL\!\big(p_\phi(z\mid x,q)\,\|\,r_\psi(z\mid q)\big).
\label{eq:a1-elbo}
\end{equation}
Minimising the empirical negative bound gives the $\qvib$ loss
\begin{equation}
\mathcal{L}_{\qvib}(\phi,\theta,\psi)
= \frac{1}{N}\sum_{i=1}^{N}\Big[
\E_{\epsilon\sim\mathcal{N}(0,I_d)}\!\big[-\log q_\theta(y_i\mid z_i)\big]
+ \beta\,\KL\!\big(p_\phi(z\mid x_i,q_i)\,\|\,r_\psi(z\mid q_i)\big)\Big],
\label{eq:a1-loss}
\end{equation}
with $z_i = \mu_\phi(x_i,q_i) + \sigma_\phi(x_i,q_i)\odot\epsilon$ (the reparameterisation trick of \citealp{kingma2014auto}).
\end{proposition}

\begin{proof}
The argument is the standard variational construction generalised to the quality-conditional
prior. It proceeds in five steps.

\paragraph{Step 1 (Importance density).}
By the law of total probability and $\int p_\phi(z\mid x,q)\,dz=1$,
\begin{equation}
\log p(y\mid x,q) = \log\!\int p_\phi(z\mid x,q)\,
\frac{p(y\mid z)\,p(z\mid x,q)}{p_\phi(z\mid x,q)}\,dz.
\label{eq:a1-p1-1}
\end{equation}

\paragraph{Step 2 (Jensen).}
The logarithm is concave, so
\begin{equation}
\log p(y\mid x,q) \;\geq\;
\E_{p_\phi(z\mid x,q)}\!\Big[\log\frac{p(y\mid z)\,p(z\mid x,q)}{p_\phi(z\mid x,q)}\Big].
\label{eq:a1-p1-2}
\end{equation}

\paragraph{Step 3 (Decompose the integrand).}
Insert the prior $r_\psi(z\mid q)$:
\begin{equation}
\log\frac{p(y\mid z)\,p(z\mid x,q)}{p_\phi(z\mid x,q)}
= \log p(y\mid z) + \log\frac{r_\psi(z\mid q)}{p_\phi(z\mid x,q)}
+ \log\frac{p(z\mid x,q)}{r_\psi(z\mid q)}.
\label{eq:a1-p1-3}
\end{equation}

\paragraph{Step 4 (Variational approximations).}
Replace the intractable likelihood $p(y\mid z)$ by $q_\theta(y\mid z)$, and use the KL identity
\begin{equation}
\E_{p_\phi}\!\Big[\log\tfrac{p(z\mid x,q)}{r_\psi(z\mid q)}\Big]
= \KL\!\big(p_\phi\,\|\,r_\psi\big) - \KL\!\big(p_\phi\,\|\,p(\cdot\mid x,q)\big).
\label{eq:a1-p1-4}
\end{equation}

\paragraph{Step 5 (Drop the non-negative residual).}
Since $\KL(p_\phi\,\|\,p(\cdot\mid x,q))\geq 0$, substituting \eqref{eq:a1-p1-3}--\eqref{eq:a1-p1-4}
into \eqref{eq:a1-p1-2} yields
\begin{equation}
\log p(y\mid x,q) \;\geq\;
\E_{p_\phi}\!\big[\log q_\theta(y\mid z)\big] - \KL\!\big(p_\phi\,\|\,r_\psi\big),
\label{eq:a1-p1-5}
\end{equation}
and introducing the trade-off multiplier $\beta$ gives \eqref{eq:a1-elbo}. The empirical form
\eqref{eq:a1-loss} follows by Monte-Carlo estimation of the expectation under the
reparameterisation trick.
\end{proof}

\begin{remark}[Strict generalisation of VIB]\label{rem:a1-vib}
When $r_\psi(z\mid q)=\mathcal{N}(0,I_d)$ for all $q$ and $p_\phi$ ignores $q$, \eqref{eq:a1-loss}
reduces \emph{exactly} to the variational information bottleneck of \citet{alemi2017deep}. Thus
$\qvib$ strictly generalises VIB by conditioning the marginal prior on perceptual quality; it is
an information-theoretic generalisation, not a posterior-reweighting scheme.
\end{remark}

\begin{remark}[Bound tightness]\label{rem:a1-tight}
The ELBO gap equals $\KL(p_\phi(z\mid x,q)\,\|\,p(z\mid x,q))$. A prior $r_\psi(z\mid q)$ matched to
$p(z\mid q)$ tightens the bound relative to a quality-agnostic $\mathcal{N}(0,I_d)$, most visibly for
low-quality inputs where $p(z\mid q)$ departs furthest from the standard Gaussian.
\end{remark}

% ═══════════════════════════════════════════════════════════════════════════════
\subsection{Quality-Adaptive Prior and Lemma 1 (Monotonicity)}\label{app:a1-lemma1}

We model the prior as an isotropic Gaussian whose variance is a smooth function of $\qbar$:
\begin{equation}
r_\psi(z\mid q) = \mathcal{N}\!\big(z;\,0,\,\sigma^2(\qbar)\,I_d\big),\qquad
\sigma^2(\qbar) = \sigma_0^2 + (1-\sigma_0^2)\,\sigmoid\!\big(-\alpha(\qbar-\tau)\big),
\label{eq:a1-prior}
\end{equation}
with baseline variance $\sigma_0^2\in(0,1]$, threshold $\tau\in[0,1]$, and sharpness $\alpha>0$.

\begin{lemma}[restated, monotonicity of $\sigma^2$]\label{lem:mono-full}
$\sigma^2(\qbar)$ is strictly decreasing in $\qbar$ on $(0,1)$.
\end{lemma}

\begin{proof}
Let $s(\cdot)=\sigmoid(\cdot)$, so $s'(\cdot)=s(\cdot)(1-s(\cdot))$. Differentiating
\eqref{eq:a1-prior} with $x=-\alpha(\qbar-\tau)$,
\begin{equation}
\frac{d\sigma^2}{d\qbar}
= (1-\sigma_0^2)\,s'\!\big(-\alpha(\qbar-\tau)\big)\cdot(-\alpha)
= -(1-\sigma_0^2)\,\alpha\,s(x)\,(1-s(x)) < 0,
\label{eq:a1-mono}
\end{equation}
since $\alpha>0$, $s(x)\in(0,1)$, and $1-\sigma_0^2\geq 0$.
\end{proof}

\paragraph{Asymptotics.}
As $\qbar\to 1$ with $\alpha(1-\tau)\to\infty$, $\sigma^2\to\sigma_0^2$ (concentrated prior, bottleneck
relaxed). As $\qbar\to 0$ with $\alpha\tau\to\infty$, $\sigma^2\to 1$ (the maximally uninformative
$\mathcal{N}(0,I_d)$, bottleneck maximally constrictive).

\paragraph{Explicit KL.}
For diagonal encoder $p_\phi=\mathcal{N}(\mu,\operatorname{diag}(\sigma_1^2,\dots,\sigma_d^2))$ and prior
\eqref{eq:a1-prior},
\begin{equation}
\KL(p_\phi\,\|\,r_\psi)
= \frac{1}{2}\sum_{j=1}^{d}\Big[\frac{\mu_j^2+\sigma_j^2}{\sigma^2(\qbar)} - 1
- \log\frac{\sigma_j^2}{\sigma^2(\qbar)}\Big].
\label{eq:a1-kl}
\end{equation}
As $\qbar$ falls, $\sigma^2(\qbar)$ rises (Lemma~\ref{lem:mono-full}); to keep \eqref{eq:a1-kl}
finite the encoder shrinks $|\mu_j|\to 0$ and inflates $\sigma_j^2\to\sigma^2(\qbar)$, pushing the
code toward the uninformative prior --- a soft, differentiable quality gate.

% ═══════════════════════════════════════════════════════════════════════════════
\subsection{Lemma 2 (Softmax Perturbation Stability)}\label{app:a1-lemma2}

\begin{lemma}[softmax $\ell_\infty\!\to\!\ell_1$ stability]\label{lem:softmax}
Let $a=\softmax(s)$ and $a'=\softmax(s+\delta)$ for $s,\delta\in\mathbb{R}^n$. Then
\begin{equation}
\|a'-a\|_1 \;\leq\; 2\max_{1\leq j\leq n}|\delta_j|.
\label{eq:a1-softmax}
\end{equation}
\end{lemma}

\begin{proof}
The softmax $f:\mathbb{R}^n\to\Delta^{n-1}$ has Jacobian $J_{jk}=f_j(\mathbb{1}[j=k]-f_k)$. For any
direction $v$, $(Jv)_j = f_j(v_j-\bar{v})$ with $\bar{v}=\sum_k f_k v_k$, so
\begin{equation}
\|Jv\|_1 = \sum_{j=1}^n f_j\,|v_j-\bar{v}|
\leq \sum_{j=1}^n f_j\big(|v_j|+|\bar{v}|\big)
\leq 2\|v\|_\infty,
\label{eq:a1-jac}
\end{equation}
where $|\bar{v}|\leq\max_k|v_k|=\|v\|_\infty$ and $\sum_j f_j=1$. Hence the induced operator norm
satisfies $\|J\|_{\infty\to 1}\leq 2$. Writing $g(t)=\softmax(s+t\delta)$ and applying the mean
value theorem,
\begin{equation}
\|a'-a\|_1 = \Big\|\int_0^1 J(s+t\delta)\,\delta\,dt\Big\|_1
\leq \sup_{t}\|J(s+t\delta)\|_{\infty\to 1}\,\|\delta\|_\infty
\leq 2\max_j|\delta_j|. \qedhere
\label{eq:a1-mvt}
\end{equation}
\end{proof}

\begin{remark}
The constant $2$ is conservative; \citet{gao2017properties} sharpen it to $1$ in certain regimes. We
retain $2$ for a clean worst-case statement in Theorem~\ref{thm:drift-full}.
\end{remark}

% ═══════════════════════════════════════════════════════════════════════════════
\subsection{Theorem 1 (Quality-Aware Attention Drift Bound)}\label{app:a1-thm1}

We inject quality into the encoder's self-attention via two learnable embeddings
$u,v:[0,1]^5\to\mathbb{R}^m$, parameterised as $u(q)=\tilde{u}(\mathbf{1}-q)$,
$v(q)=\tilde{v}(\mathbf{1}-q)$ with $\tilde{u}(0)=\tilde{v}(0)=0$, so that perfect quality
($q=\mathbf{1}$) recovers standard attention. The modulated logits are
$\tilde{s}_{ij}=x_i^\top x_j/\sqrt{d}+\delta_{ij}$, $\delta_{ij}=u(q_i)^\top v(q_j)$.

\begin{theorem}[restated, attention drift bound]\label{thm:drift-full}
Let $\tilde{u},\tilde{v}$ be $L_{\tilde{u}}$- and $L_{\tilde{v}}$-Lipschitz in $\ell_2$ with
$\tilde{u}(0)=\tilde{v}(0)=0$, and let $\varepsilon=\max_i\|\mathbf{1}-q_i\|_\infty$ be the
worst-case quality deviation. With $a_i,a'_i$ the attention rows before/after modulation,
\begin{equation}
\max_{1\leq i\leq n}\|a'_i-a_i\|_1 \;\leq\; 10\,L_{\tilde{u}}L_{\tilde{v}}\,\varepsilon^2.
\label{eq:a1-drift}
\end{equation}
\end{theorem}

\begin{proof}
Four steps.

\paragraph{Step 1 (Cauchy--Schwarz on the bias).}
$|\delta_{ij}| = |\tilde{u}(\mathbf{1}-q_i)^\top\tilde{v}(\mathbf{1}-q_j)|
\leq \|\tilde{u}(\mathbf{1}-q_i)\|_2\,\|\tilde{v}(\mathbf{1}-q_j)\|_2.$

\paragraph{Step 2 (Lipschitz with zero anchor).}
Since $\tilde{u}(0)=0$, $\|\tilde{u}(\mathbf{1}-q_i)\|_2\leq L_{\tilde{u}}\|\mathbf{1}-q_i\|_2$, and
likewise for $\tilde{v}$. Hence
\begin{equation}
|\delta_{ij}| \leq L_{\tilde{u}}L_{\tilde{v}}\,\|\mathbf{1}-q_i\|_2\,\|\mathbf{1}-q_j\|_2.
\label{eq:a1-bias}
\end{equation}

\paragraph{Step 3 (Apply Lemma~\ref{lem:softmax}).}
For each row $i$, $\|a'_i-a_i\|_1\leq 2\max_j|\delta_{ij}|$, so
\begin{equation}
\max_i\|a'_i-a_i\|_1
\leq 2L_{\tilde{u}}L_{\tilde{v}}\max_{i,j}\big(\|\mathbf{1}-q_i\|_2\,\|\mathbf{1}-q_j\|_2\big).
\label{eq:a1-row}
\end{equation}

\paragraph{Step 4 ($\ell_2\to\ell_\infty$ on a $5$-vector).}
For $q\in[0,1]^5$, $\|\mathbf{1}-q\|_2^2=\sum_{d=1}^5(1-q_d)^2\leq 5\,\varepsilon^2$, i.e.\
$\|\mathbf{1}-q\|_2\leq\sqrt{5}\,\varepsilon$. Substituting into \eqref{eq:a1-row},
\begin{equation}
\max_i\|a'_i-a_i\|_1 \leq 2L_{\tilde{u}}L_{\tilde{v}}\,(\sqrt{5}\,\varepsilon)^2
= 10\,L_{\tilde{u}}L_{\tilde{v}}\,\varepsilon^2. \qedhere
\end{equation}
\end{proof}

\begin{corollary}[Spectral-norm control]\label{cor:a1-spectral}
If $\tilde{u},\tilde{v}$ are spectrally normalised \citep{miyato2018spectral}
($L_{\tilde{u}}=L_{\tilde{v}}=1$), the drift is at
most $10\varepsilon^2$; for a moderately degraded input ($\varepsilon=0.3$) the attention rows shift by
at most $0.9$ in $\ell_1$ --- a bounded, not arbitrary, reorganisation.
\end{corollary}

\begin{remark}[Near-tightness]\label{rem:a1-drift-tight}
With scalar embeddings $\tilde{u}(x)=\tilde{v}(x)=Lx$ ($m=1$), a two-patch input at $q=\mathbf{0}$
($\varepsilon=1$) drives the drift arbitrarily close to $2L^2$, against the bound $10L^2$; the factor
$5$ gap stems from the generic $\ell_2\to\ell_\infty$ conversion and the constant in
Lemma~\ref{lem:softmax}.
\end{remark}

% ═══════════════════════════════════════════════════════════════════════════════
\subsection{Proposition 2 (Quality-Calibrated Entropy)}\label{app:a1-prop2}

Let $\hat{p}_q(y\mid x)=\E_{z\sim p_\phi(z\mid x,q)}[q_\theta(y\mid z)]$ be the predictive
distribution and $\Hent(\hat{p}_q)=-\sum_y \hat{p}_q\log\hat{p}_q$ its entropy.

\begin{proposition}[restated, entropy monotonicity]\label{prop:entmono-full}
Under the $\qvib$ objective \eqref{eq:a1-loss} with $\beta>0$, the expected predictive entropy is
non-increasing in quality:
\begin{equation}
\E_{x\sim p(x\mid\qbar_1)}\!\big[\Hent(\hat{p}_{q_1})\big]
\;\geq\; \E_{x\sim p(x\mid\qbar_2)}\!\big[\Hent(\hat{p}_{q_2})\big]
\quad\text{whenever}\quad \qbar_1<\qbar_2.
\label{eq:a1-entmono}
\end{equation}
\end{proposition}

\begin{proof}
Four steps relating $\sigma^2(\qbar)$ to the optimal encoder and thence to entropy.

\paragraph{Step 1 (Objective at fixed $\qbar$).}
The encoder $\mathcal{N}(\mu,\Sigma)$ minimises
\begin{equation}
\mathcal{J}(\mu,\Sigma) = \E_\epsilon\!\big[-\log q_\theta(y\mid\mu+\Sigma^{1/2}\epsilon)\big]
+ \beta\,\KL\!\big(\mathcal{N}(\mu,\Sigma)\,\|\,\mathcal{N}(0,\sigma^2 I_d)\big),
\label{eq:a1-p2-1}
\end{equation}
with the KL term $\tfrac{1}{2}[\|\mu\|^2/\sigma^2 + \tr(\Sigma)/\sigma^2 - d - \log\det\Sigma + d\log\sigma^2]$.

\paragraph{Step 2 (First-order condition on $\mu$).}
$\nabla_\mu\mathcal{J}=\nabla_\mu\E_\epsilon[-\log q_\theta] + (\beta/\sigma^2)\mu = 0$ gives
\begin{equation}
\mu^\star = -\frac{\sigma^2}{\beta}\,\nabla_\mu\E_\epsilon[-\log q_\theta].
\label{eq:a1-p2-2}
\end{equation}
The classification-loss gradient has bounded norm (Lipschitz for standard heads), so $\mu^\star$ is
scaled toward $0$ as $\sigma^2$ grows.

\paragraph{Step 3 (Code $\to$ uninformative prior $\Rightarrow$ max entropy).}
As $\mu\to 0$ and $\Sigma\to\sigma^2 I_d$, $z$ approaches a sample from $\mathcal{N}(0,\sigma^2 I_d)$ and
$\hat{p}_q(y\mid x)\to\E_{z\sim\mathcal{N}(0,\sigma^2 I_d)}[q_\theta(y\mid z)]$, which carries no
discriminative signal and tends to the uniform distribution (entropy $\log K$, the maximum).

\paragraph{Step 4 (Monotonicity).}
By Lemma~\ref{lem:mono-full}, $\qbar_1<\qbar_2\Rightarrow\sigma^2(\qbar_1)>\sigma^2(\qbar_2)
\Rightarrow\|\mu^\star(\qbar_1)\|\leq\|\mu^\star(\qbar_2)\|$ (Step~2)
$\Rightarrow\Hent(\hat{p}_{q_1})\geq\Hent(\hat{p}_{q_2})$ (Step~3). The last implication is the data
processing inequality on $Y\to X\to Z\to\hat{Y}$: a more compressed $Z$ cannot raise predictive
information.
\end{proof}

\begin{remark}[Calibration without post-hoc tuning]\label{rem:a1-calib}
Proposition~\ref{prop:entmono-full} certifies that $\qvib$ is provably more uncertain on lower-quality
inputs, structurally and end-to-end --- in contrast to hard quality gating (non-differentiable,
binary accept/reject) and post-hoc temperature scaling (applied after training).
\end{remark}

\paragraph{Scope of the guarantee.}
Steps 2--3 hold at the training equilibrium of \eqref{eq:a1-loss} for classifiers with bounded
gradient and a well-fit $q_\theta$; \eqref{eq:a1-entmono} is therefore a guarantee about the
optimiser's fixed point, not a finite-sample certificate. The empirical monotone trend
($\rho(\Hent,\qbar)<0$) is reported in the main results.
        % TODO: Prop 1 / Lemma 1 / Thm 1 / Prop 2
% % Appendix A2.3 — full proof of Theorem 2 (closed-loop agent expected-risk bound).
% Source-of-truth: project/plans/Theorem2_agent_risk_bound.md (audited 2026-05-27).
% Numerical toy: tests/test_theorems_numerical.py (P1/P2/P3 + Cor 2.1).

\subsection{Proof for the Closed-Loop Agent (Theorem 2)}\label{app:a2-thm2}

\subsubsection{Decision-theoretic setup}\label{app:a2-thm2-setup}
At input $x$ with quality scalar $\qbar(x)$ already scored by \visiscore, the agent
selects one of four actions $\mathcal{A}=\{\actd,\acte,\actq,\actr\}$ (direct, enhance, query,
refuse) with constant costs $c_d=0,\,c_e,\,c_q,\,c_r\geq 0$. The per-action expected risk is
\begin{equation}
\mathcal{R}_a(x) := \E_{y\sim p(y\mid x)}\!\big[\ell(\hat{y}_a(x), y)\big] + c_a,
\qquad a\in\mathcal{A},
\label{eq:per-action}
\end{equation}
with bounded loss $\ell\in[0,M]$ ($M=1$ for $0$--$1$ loss). For the enhance action the
expectation is still taken over $p(y\mid x)$ (the true label is fixed by the original image),
while the predictor acts on $\Tw(x,q)$ --- the diagnosis-preserving setting.

The threshold policy $\pi^{*}$ uses $\qbar$ and a secondary entropy gate $\hhigh$:
\begin{equation}
\pi^{*}(x) =
\begin{cases}
\actr & \qbar < \taulow \;\text{ or }\; \big(\qbar\in[\tauenh,\tauhigh]\text{ and }\Hent(\hat{p}_x)>\hhigh\big)\\
\actq & \taulow \leq \qbar < \tauenh \\
\acte & \tauenh \leq \qbar < \tauhigh \text{ and } \Hent(\hat{p}_x)\leq \hhigh \\
\actd & \qbar \geq \tauhigh
\end{cases}
\label{eq:agent-policy}
\end{equation}
with $0<\taulow<\tauenh<\tauhigh<1$ chosen by validation grid search; $\hhigh$ (the $75$th
percentile of the validation entropy distribution) separates ``quality-sufficient but
intrinsically ambiguous'' cases that enhancement cannot salvage.

\begin{assumption}[Unimodal enhancement gain]\label{as:unimodal}
The expected gross enhancement gain $\qbar\mapsto\E[\Gen(x)\mid\qbar(x)=\qbar]$ is unimodal in
$\qbar$ on $(0,1)$. This is an empirical modelling assumption (supported by the salvage-rate
curve), not a derived property; the threshold construction in Lemma~\ref{lem:window} relies on
it to guarantee that the effective enhancement band is a single connected interval.
\end{assumption}

\subsubsection{Auxiliary lemmas}\label{app:a2-thm2-lemmas}
\begin{lemma}[Entropy--risk coupling]\label{lem:entropy-risk}
For a calibrated predictor with TV calibration error $\delta_{\TV}$ and bounded
$\ell\in[0,M]$,
$\Rdirect(x) \leq M\big(1-\exp(-\Hent(\hat{p}_x))\big) + M\,\delta_{\TV}(\hat{p},p)$.
\end{lemma}
\begin{proof}
The optimal plug-in classifier has $0$--$1$ risk $1-\max_y\hat{p}(y\mid x)\leq 1-\exp(-\Hent(\hat{p}_x))$
by Gibbs' inequality $\max_y p(y)\geq\exp(-\Hent(p))$. The calibration term follows from
$|\E_p[\ell]-\E_{\hat{p}}[\ell]|\leq M\,\delta_{\TV}(\hat{p},p)$.
\end{proof}

\begin{lemma}[Enhancement risk reduction]\label{lem:enh-risk}
Define the gross gain through the post-enhancement predictive entropy,
$\Gen(x):=\Rdirect(x)-\big(1-e^{-\E[\Hent_{\Tw(x,q)}]}\big)$. Then, using only
Lemma~\ref{lem:entropy-risk} (entropy--risk coupling) and the DP-loss calibration bound of
Lemma~\ref{lem:dp-full},
\begin{equation}
\Rdirect(x) - \Renh(x) \;\geq\; \Gen(x) - c_e - M\delta_{\TV}.
\label{eq:enh-gain}
\end{equation}
On the salvage band where $\Gen(x)>0$ (an empirical condition, equivalently a positive
$\Delta\mathrm{AUC}_{\mathrm{enh}}$; validated by experiment E7), the right-hand side is
positive whenever $\Gen(x)>c_e+M\delta_{\TV}$.
\end{lemma}
\begin{proof}[Derivation (under the stated assumptions; see Threats to validity)]
Apply Lemma~\ref{lem:entropy-risk} to the enhanced input:
$\Renh(x)-c_e \leq 1-\exp(-\E[\Hent_{\Tw(x,q)}]) + M\delta_{\TV}$, then rearrange against
$\Rdirect(x)$ to obtain \eqref{eq:enh-gain}. We do \emph{not} invoke a net predictive-entropy
reduction from Proposition~\ref{prop:entropy-full} (which is not empirically supported);
instead, $\Gen(x)$ is defined directly from the post-enhancement entropy, and its positivity
on the salvage band is an empirical premise carried by the DP-loss mutual-information lower
bound (Lemma~\ref{lem:dp-full}, experiment E7) rather than a consequence of
Proposition~\ref{prop:entropy-full}. Accordingly $\Gen(x)$ is not asserted to be non-negative
for every $x$; it is positive only on the empirically identified band.
\end{proof}

\begin{lemma}[Enhancement effectiveness window]\label{lem:window}
There exist $\tauenh<\tauhigh$ such that $\Gen(x) > c_e + M\delta_{\TV}$ for almost all $x$ with
$\qbar(x)\in(\tauenh,\tauhigh)$, and $\Gen(x)\leq c_e+M\delta_{\TV}$ outside.
\end{lemma}
\begin{proof}[Proof sketch]
As $\qbar\to 1$ the predictor approaches a one-hot, $\Hent\to 0$ and $\Gen\to 0$, giving an
upper threshold $\tauhigh$ by continuity. For $\qbar<\tauenh$ (severe degradation, e.g.\ loss
of completeness $q_3$), $\Tw$ cannot raise quality, so $\E[\Hent_{\mathrm{enh}}]\approx\E[\Hent_{\mathrm{direct}}]$
and $\Gen\approx 0$. Inside the band, Proposition~\ref{prop:entropy-full} gives a significant
entropy drop, hence $\Gen>c_e$. The thresholds are
$\tauenh=\inf\{\qbar:\E[\Gen]>c_e+M\delta_{\TV}\}$ and
$\tauhigh=\sup\{\qbar:\E[\Gen]>c_e+M\delta_{\TV}\}$; under
Assumption~\ref{as:unimodal} ($\E[\Gen]$ unimodal in $\qbar$), the super-level set
$\{\qbar:\E[\Gen]>c_e+M\delta_{\TV}\}$ is a single connected interval, so the band is a
non-empty connected set.
\end{proof}

\begin{lemma}[Query/refuse safety]\label{lem:query-refuse}
For $\qbar<\tauenh$: (i) if $\qbar\geq\taulow$ and the expected retake quality exceeds
$\tauhigh$, then $\mathcal{R}_{\actq}(x)\leq\Rdirect(x)-(\Gret(x)-c_q)$; (ii) if $\qbar<\taulow$,
then $\mathcal{R}_{\actr}(x)=c_r+\ell_r\leq\Rdirect(x)$ in expectation over $x$.
\end{lemma}
\begin{proof}[Proof sketch]
Identical to Lemma~\ref{lem:enh-risk} with $\Tw$ replaced by the stochastic retake operator
$P_{\mathrm{retake}}$ (satisfying $\E[\Hent_{\mathrm{retake}}]\leq\Hent(\hat{p}_x)$) and $c_e$
by $c_q$. The refuse branch is a constant referral cost.
\end{proof}

\subsubsection{Main result}\label{app:a2-thm2-main}
\begin{theorem}[Closed-loop agent expected-risk bound]\label{thm:agent}
Let $\Tw$ be a diagnosis-preserving enhancement satisfying Lemma~\ref{lem:dp-full} with
calibration error $\delta_{\TV}$ (its information-preservation guarantee; we do not rely on a
net-entropy claim from Proposition~\ref{prop:entropy-full}), let
Assumption~\ref{as:unimodal} hold, and let $\pi^{*}$ be the threshold policy
\eqref{eq:agent-policy}. Then for any $x$,
\begin{equation}
\Ragent(x) \;\leq\; \Rdirect(x) - \Delta(\qbar(x),\Tw),
\qquad
\Delta(\qbar,\Tw) > 0 \;\text{ if }\; \qbar\in(\tauenh,\tauhigh)
\text{ and } \E[\Gen]>c_e+M\delta_{\TV},
\label{eq:thm2}
\end{equation}
and on the enhance channel
$\Delta(\qbar,\Tw)\big|_{\qbar\in(\tauenh,\tauhigh)}
= \E_{x:\qbar(x)=\qbar}[\Gen(x)] - c_e - M\delta_{\TV}$.
Outside this sub-band the policy reverts to query/refuse/direct and $\Delta=0$, so the bound
holds with equality; we state the strict improvement as a sufficient condition on the
enhancement band rather than as an ``if and only if''.
\end{theorem}
\begin{proof}
We consider the four cases of the policy partition \eqref{eq:agent-policy}.

\emph{Case 1 ($\qbar\geq\tauhigh$, direct).} $\Ragent(x)=\Rdirect(x)$; by Lemma~\ref{lem:window}
$\Gen\leq c_e+M\delta_{\TV}$ so $\Delta=0$ and \eqref{eq:thm2} holds with equality.

\emph{Case 2 ($\qbar\in[\tauenh,\tauhigh)$ and $\Hent(\hat{p}_x)\leq\hhigh$, enhance).}
$\Ragent(x)=\Renh(x)$. Taking the conditional expectation of Lemma~\ref{lem:enh-risk} at
$\qbar(x)=\qbar$,
$\E[\Ragent]\leq\E[\Rdirect]-(\E[\Gen]-c_e-M\delta_{\TV})$, and by Lemma~\ref{lem:window} the
parenthesised term is strictly positive, giving $\Delta>0$.

\emph{Case 3 ($\qbar\in[\taulow,\tauenh)$, query).} By Lemma~\ref{lem:query-refuse}(i),
$\Ragent(x)\leq\Rdirect(x)-(\Gret(x)-c_q)$. On this band $\Gen\leq c_e+M\delta_{\TV}$, so the
enhance-channel gain $\Delta:=0$, consistent with \eqref{eq:thm2}.

\emph{Case 4 ($\qbar<\taulow$, refuse).} By Lemma~\ref{lem:query-refuse}(ii),
$\Ragent(x)=c_r+\ell_r\leq\Rdirect(x)$ in expectation; $\Delta=0$.

The four cases establish \eqref{eq:thm2}.
\end{proof}

\begin{corollary}[Agent never worse than direct]\label{cor:never-worse}
Under the thresholds of Lemmas~\ref{lem:window}--\ref{lem:query-refuse},
$\Ragent(x)\leq\Rdirect(x)$ for all $x$.
\end{corollary}

\begin{corollary}[Strict improvement on the salvage band]\label{cor:population}
If $\pisalv := \mathbb{P}[\qbar(X)\in[\tauenh,\tauhigh]] > 0$, then
$\E_x[\Ragent(X)] \leq \E_x[\Rdirect(X)] - \pisalv\cdot\Delta_{\mathrm{avg}}$ with
$\Delta_{\mathrm{avg}}>0$.
\end{corollary}

\paragraph{Tight constants.}
With the \qvib{} predictor the calibration proxy is $\delta_{\TV}\approx\ECE=0.098$; the
normalised enhancement cost is $c_e=0.02$; a validation grid yields
$\tauenh\approx 0.35$, $\tauhigh\approx 0.55$. Corollary~\ref{cor:population} maps directly to
the empirical salvage-rate ($\pisalv$) and $|\Delta\mathrm{AUC}|$ metrics.

\paragraph{Threats to validity.}
The Gibbs lower bound in Lemma~\ref{lem:entropy-risk} is tighter for $K=2$ than for large $K$;
the calibration term treats ECE as a TV proxy (a piecewise-linear calibration map closes the
gap, Appendix~A2.3.5); the unimodality of $\Gen$ (Assumption~\ref{as:unimodal}) is an empirical
modelling assumption rather than a derived property, and the salvage-band non-emptiness it
underwrites holds only when $\Tw$ does not collapse (Stage-1 PSNR$\,\geq 30$\,dB gate); the
positivity of the gross gain $\Gen$ rests on the DP-loss mutual-information lower bound
(Lemma~\ref{lem:dp-full}, experiment E7), not on a net predictive-entropy reduction from
Proposition~\ref{prop:entropy-full}, which is not empirically supported; and
Theorem~\ref{thm:agent} is an expected-risk bound, with a Hoeffding finite-sample companion
$|\hat{\Delta}_{\mathrm{avg}}-\Delta_{\mathrm{avg}}|\leq M\sqrt{\log(2/\delta)/(2n)}$
(Appendix~A2.3.6).
  % TODO: Theorem 2 full proof
% % Appendix A3 — full proof of Corollary 1 (composite quality-conditional calibration ECE bound).
% Source-of-truth: project/plans/Corollary1_qvib_qcts_ece_bound.md (audited 2026-05-27).
% QCTS = quality-conditional temperature scaling from a prior anonymous submission
% (cited upon acceptance); referred to via the \qcts macro.

\section{Proof of the Composite Calibration Bound (Corollary 1)}\label{app:a3}

Let $\pQV$ be the \qvib{} predictive (Proposition~1) and $T(\qbar)=\log(1+\exp(T_0+\alpha(1-\qbar)))$
the \qcts{} temperature schedule with Lipschitz constant $\Lt$ and minimum $\Tmin$. The
composite $\pComp$ rescales the \qvib{} logits by $T(\qbar)$.

\begin{corollary}[Composite calibration ECE bound]\label{cor:composite}
\begin{equation}
\ECE(\pComp) \;\leq\; \ECE(\pQCTS) + \epsqts,
\qquad
\epsqts = O\!\Big(\tfrac{\Lt\,\|\ell\|_{\max}}{\Tmin^2}\cdot\sigma_{\qbar\mid c}\Big),
\label{eq:cor1}
\end{equation}
where $\epsqts$ vanishes when $T\equiv 1$ (composite $=\pQV$) or
$\Var[\qbar\mid\pQV]=0$ (quality-deterministic output). We state the one-sided bound against
$\ECE(\pQCTS)$ only: the stronger two-sided form
$\min(\ECE(\pQV),\ECE(\pQCTS))+\epsqts$ would additionally require the $\qbar$-averaged
conditional reliability gap to be no larger than \emph{either} marginal calibration error,
which we do not establish here (the $\qbar$-average can exceed the better of the two marginals;
see the remark after the proof).
\end{corollary}

\begin{proof}[Derivation (under the stated Lipschitz schedule assumptions; see the main-text
Limitations for non-watertight regimes)]
The argument proceeds in four steps and establishes the one-sided bound \eqref{eq:cor1}.

\paragraph{Step 1 (ECE as a reliability integral).}
By the Brier/Murphy decomposition, with $c=\hat{p}(Y=1\mid X)$, calibration curve
$\pi(c)=\E[Y\mid\hat{p}=c]$ and density $\rho_{\hat{p}}$,
\begin{equation}
\ECE(\hat{p}) = \int_0^1 |c - \pi(c)|\,\rho_{\hat{p}}(c)\,dc.
\label{eq:cor1-1}
\end{equation}

\paragraph{Step 2 (effect of \qcts{} rescaling).}
With $\ell(x)=\mathrm{logits}(\pQV(\cdot\mid x))$, the composite marginal satisfies
$c^{(\mathrm{comp})} = \sigmoid(\mathrm{logit}(c^{(\mathrm{QV})})/T(\qbar))$, a
$\qbar$-dependent monotone bijection of $[0,1]$ for each fixed $\qbar$.

\paragraph{Step 3 (Lipschitz drift of the calibration curve).}
The composite curve averages the $\qbar$-conditional curves,
$\pi^{(\mathrm{comp})}(c)=\E_{\qbar\mid\pComp=c}[\pi(c;\qbar)]$. Since $T$ is $\Lt$-Lipschitz in
$\qbar$, the chain rule through softmax/sigmoid gives
\begin{equation}
|c^{(\mathrm{comp})}(x,\qbar_1)-c^{(\mathrm{comp})}(x,\qbar_2)|
\leq \frac{\Lt}{\Tmin^2}\,|\ell(x)|\,|\qbar_1-\qbar_2|.
\label{eq:cor1-3}
\end{equation}

\paragraph{Step 4 (combine via the Murphy decomposition).}
Substituting Step~3 into \eqref{eq:cor1-1} and applying Jensen and Cauchy--Schwarz,
\begin{equation}
\ECE(\pComp)
\leq \underbrace{\E_{\qbar}\!\int|c-\pi(c;\qbar)|\rho\,dc}_{\leq\,\ECE^{(\mathrm{QCTS})}}
+ \underbrace{\tfrac{\Lt\,\|\ell\|_{\max}}{\Tmin^2}}_{=:K_T}\cdot\sigma_{\qbar\mid c},
\label{eq:cor1-4}
\end{equation}
where the first term is bounded by the calibration error of the $\qcts$-rescaled predictive
$\ECE^{(\mathrm{QCTS})}$ --- the temperature schedule $T(\qbar)$ is exactly the map whose
calibration this term measures --- and the second is the Lipschitz residual
$\epsqts := K_T\,\sigma_{\qbar\mid c}$. This gives the one-sided bound \eqref{eq:cor1}.
\end{proof}

\begin{remark}[Why the bound is one-sided]\label{rem:cor1-onesided}
We do not bound the first term of \eqref{eq:cor1-4} by $\ECE^{(\mathrm{QV})}$. The
$\qbar$-averaged conditional reliability gap $\E_{\qbar}\!\int|c-\pi(c;\qbar)|\rho\,dc$ need
not be smaller than the $\pQV$ marginal calibration error: averaging conditional absolute
deviations can exceed the absolute deviation of the marginal calibration curve (a
Jensen-type inequality runs in the unfavourable direction here). We therefore anchor the bound
on $\ECE^{(\mathrm{QCTS})}$, the calibration that the composite directly inherits through the
temperature map, and report the one-sided guarantee. % VERIFIED (会话35 theory reviewer): one-sided direction holds (q-bar-averaged conditional ECE <= ECE(QCTS), the calibration QCTS directly inherits); min(QV,QCTS) was unprovable as QV side runs Jensen the wrong way. Bound value 0.079+0.037=0.116 unchanged since min(0.098,0.079)=0.079.
\end{remark}

\paragraph{Tight constants.}
With $\alpha=0.955$ the schedule has $\Lt\approx 0.239$ and $\Tmin\approx 1.44$; the \qvib{}
logit range is $\|\ell\|_{\max}\approx 4$, giving $K_T=\Lt\|\ell\|_{\max}/\Tmin^2\approx 0.461$.
The quality spread at fixed composite confidence is $\sigma_{\qbar\mid c}\approx 0.08$, so
$\epsqts\approx 0.037$. With $\ECE(\pQCTS)=0.079$, the one-sided bound \eqref{eq:cor1} predicts
$\ECE(\pComp)\leq 0.079+0.037 = 0.116$ (for reference, the $\pQV$ marginal is
$\ECE(\pQV)=0.098$); the residual is bounded away from vacuity for any $\Lt\leq 5$ (the schedule
enforces $\alpha\leq 1$).

\paragraph{Role in the narrative.}
Corollary~\ref{cor:composite} shows the composite \emph{cannot} reach $\ECE^{(\mathrm{QCTS})}$
exactly --- there is always a Lipschitz residual $\epsqts>0$ --- so a measured composite ECE
near the prior-submission value reflects an additive $\epsqts$ rather than reuse of that number.
This makes explicit the distinction between the prior post-hoc temperature scaling and the
present end-to-end system.
  % TODO: Corollary 1 full proof

\end{document}
